% Supplemental section; no preamble.  Uses the report's theorem environments
% and the existing bibliography keys pa194, cached, and cz.
\section{PA194D: a different logical emitter with 194 states}
\label{sec:gammad}

There is a second 194-state construction, denoted
\(\mathsf{PA194D}\), which changes the numerical axiom command itself.
The separate diagonal clause in the earlier machines is
\(\Gamma(A)=A\to U(A,A)\).  PA194D supplies
\(\Gamma(D)=D\to U(D,D)\) instead, while preserving the parameter
choices in all six other heads and in the guarded separation gate.
The resulting proof system is the same numerical calculus, but its
command interface and its physical reader allocation are different.
Its state count is
\[
 134\text{ fixed controls}+69\text{ raw readers}
       -9\text{ compatible folds}=194.
\]
The command controller is the cached controller with the scheduler's
known-zero decrement, as in the supplied PA194CZ construction
\cite{cached,cz}.  In addition, the final AX return uses a decrement of
the cleared register \(D\).  The argument below separates the logical
change from these two reachable source specializations.

The dense binary table has working controls \(0,\ldots,193\), blank
start control 0, and terminal halt label 194.  Its SHA-256 is
\begin{quote}\small
\nolinkurl{4226b7702de554df4b52d10239b2bd7ad79ca90bd8faf2bd256552fe0c7fe81b}.
\end{quote}
The native-label table starts at 74, halts at 254, and has SHA-256
\begin{quote}\small
\nolinkurl{bc8ab6c58badd137535938c9810e39d19152c49dd2eb26ef4ea07dbcf3157231}.
\end{quote}
The table uses the same two-way infinite binary tape convention as
PA194J. It has 388 defined transition rows. The authoritative file is
\file{logical_emitter/machine/pa194d.tm}; its complete contents are printed
in Appendix~\ref{app:table_d}.

\subsection{The exact logical change}

Use Cantor pairing \(p\), implication \(I\), membership \(M\),
universal quantification \(U\), and equality \(E\) with the numerical
definitions in Section~\ref{sec:calibration}.  In particular,
\(U(i,B)=4p(B,i)+2\), \(\neg B=I(B,1)\), and
\(\exists x_i B=\neg U(i,\neg B)\).  Write \([a;w]=p(a,w)\),
with the leftmost entry of a displayed list at its top.  The six retained
heads, in their original push order, are
\begin{align*}
 h_1(D,H)&=(\exists x_D U(D,H))\to H,\\
 h_2(D,A,H)&=U(D,U(A,H))\to U(A,U(D,H)),\\
 h_3(D,A,H)&=M(H,D)\to(E(H,A)\to M(A,D)),\\
 h_4(D,A,H)&=E(H,D)\to(M(A,H)\to M(A,D)),\\
 h_5(D,H)&=\exists x_D E(H,D),\\
 h_6(D,A,H)&=U(D,A)\to(U(D,A\to H)\to U(D,H)).
\end{align*}
The dominated propositional axiom is
\[
 \rho_8(D,A,H,u)=((H\to u)\to D)
                         \to((D\to H)\to(A\to H)),\qquad u\le D.
\]
Let \(G(a,b,c)\) be the retained pointed-separation formula from the
supplied separate-head calculus, and put
\[
 K(a,b,c)=
 \begin{cases}G(a,b,c),&c\ne a\text{ and }c\ne b,\\
 0,&\text{otherwise}.
 \end{cases}
\]
For precision, its construction is
\begin{align*}
 V(a,b,c)&=M(a,b)\to U(c,M(a,b)),\\
 P(a,b,c)&=c\to M(c,b),\\
 S_*(a,b,c)&=U\bigl(b,(M(a,b)\to P(a,b,c))
                                \wedge(M(c,b)\to c)\bigr),\\
 Q_*(a,b,c)&=U(c,S_*(a,b,c)\to\neg P(a,b,c)),\\
 G(a,b,c)&=\neg(V(a,b,c)\to Q_*(a,b,c)),
 \qquad B\wedge C=\neg(B\to\neg C).
\end{align*}

For an incoming store \(\langle t,w,d,h\rangle\), define
\(A=c\), \(D=d+x\), and \(H=h+y\).  The new command is exactly
\[
\begin{aligned}
 \mathsf{AX}^{D}_{c,x,y}\langle t,w,d,h\rangle
   &=\langle K(H,D,A),W',0,0\rangle,\\
 W'&=[\Gamma(D),h_6,h_5,h_4,h_3,h_2,h_1,\\[-2pt]
    &\hspace{3.5em}\rho_8(D,A,H,d),U(d,h),t;w].
\end{aligned}
\tag{D-AX}\label{eq:gammad-ax}
\]
Every head uses the displayed \(D,A,H\); only the argument of the
separate \(\Gamma\) has changed.  Thus the command still pushes ten
entries, counting the saved top, and still clears \(d,h\).
A failed gate guard does not remove the separately pushed
\(\Gamma(D)\).

This change is visible even on a zero incoming store.  For the command
parameters \((c,x,y)=(0,1,0)\), the former first stack entry is
\(\Gamma(0)=20\), whereas the new one is \(\Gamma(1)=832\).
The other emitted entries retain their former values.  In particular,
this is a change of the correlation between axiom parameters, not a
global exchange of physical register names.

Let \(\mathcal C_s\) be the supplied separate-head calculus: the six
families above, dominated \(\rho_8\), every \(\Gamma(B)\), each
guarded \(G(a,b,c)\), modus ponens, and theorem generalization
\(B\mapsto U(i,B)\).  ``Theorem'' means that there are no undischarged
hypotheses; the formula may have free variable labels.  The inherited
calibration gives primitive recursive proof translations, with
correctness provable in PA, establishing
\[
 \mathrm{PA}\vdash
 \left(\exists q\,\operatorname{Prf}_{\mathcal C_s}(q,1)
                   \ \leftrightarrow\ \neg\operatorname{Con}(\mathrm{PA})\right).
 \tag{D-Cal}\label{eq:gammad-cal}
\]
This is the separate-head calibration retained in
\file{references/logical_calibration/separate_head_calculus.tex},
together with its diagonal-calculus and PA195 dependencies
\cite{pa194}.  PA194D uses precisely this calculus.  It requires neither
a new logical axiom nor a new translation of PA proofs.

\subsection{All-store command identities and the changed selector}

For this subsection, use the ordinary cached-C source with the new
emitter and with its two original unconditional jumps.  Its numerical
commands have arbitrary stores \(\langle t,w,d,h\rangle\); scratch
registers obey the inherited clean command-boundary convention.  Write
\(S\) for cached Shift:
\[
 S\langle t,w,d,h\rangle=
 \begin{cases}
 \mathsf{halt},&h=1,\\
 \langle h,\operatorname{tl}(w),d,\operatorname{hd}(w)\rangle,&h\ne1.
 \end{cases}
\]
For \(y=d+x\), cached modus ponens is
\[
 \mathsf{MP}_x\langle t,w,d,h\rangle=
 \begin{cases}
 \langle0,w,0,y\rangle,&I(t,y)=h,\\
 \langle0,w,y,0\rangle,&I(t,y)\ne h.
 \end{cases}
\]
These are exactly the commands proved for the supplied cached
controller \cite{cached}.

Number the generated stack entries in (\ref{eq:gammad-ax}) by
\(q_1,\ldots,q_9\), so \(q_1=\Gamma(D)\),
\(q_8=\rho_8(D,A,H,d)\), and \(q_9=U(d,h)\).
The first eight entries and the gate have residue zero modulo four;
\(q_9\) has residue two.  None is code one.  Consequently the
ten-entry stack calculations in the supplied cached command audit apply
without any change in their exponents.  Explicitly, define
\begin{align*}
 \mathsf{Push}&=\mathsf{AX}^{D}_{0,0,0};S^9;\mathsf{MP}_0,\\
 \mathsf{Pop}&=\mathsf{AX}^{D}_{0,0,0};S^{11},\\
 \mathsf{Discard}&=\mathsf{MP}_1;\mathsf{AX}^{D}_{1,0,0};S^{10};\mathsf{MP}_0.
\end{align*}
Their exact contracts are
\begin{align*}
 \mathsf{Push}\langle t,w,d,h\rangle
   &=\langle0,[t;w],0,0\rangle,\\
 \mathsf{Pop}\langle t,[a;w],d,h\rangle
   &=\begin{cases}\mathsf{halt},&t=1,\\
                    \langle t,w,0,a\rangle,&t\ne1,
      \end{cases}\\
 \mathsf{Discard}\langle t,w,0,0\rangle
   &=\langle0,w,0,0\rangle.
\end{align*}
For example, \(\mathsf{AX}^{D};S^9\) leaves
\(\langle q_8,[t;w],0,q_9\rangle\), and
\(\mathsf{AX}^{D};S^{10}\) leaves
\(\langle q_9,w,0,t\rangle\).  The first ten Shifts test only zero
and generated entries; the eleventh tests the saved top.  Push's last
MP fails because an implication has residue zero and \(q_9\) has
residue two.  Discard's first MP fails since \(I(t,1)>0\); after
its AX and Shifts, the last MP fails since \(I(U(1,0),0)>0\).
These proofs use the total projection equations
\(\operatorname{hd}[a;w]=a\) and
\(\operatorname{tl}[a;w]=w\), including the alias \([0;0]=0\).
They do not presume that a zero push increases a stored stack length.

For \(1\le j\le7\), define a selector whose argument order is
\emph{\(D,H,A\)}:
\[
 \mathsf{select}_{j}(D,H,A)=
 \mathsf{Discard};\mathsf{AX}^{D}_{A,D,H};S^{j+1};
                                      \mathsf{Pop}^{9-j}.
\]
It sends every \(\langle t,w,0,0\rangle\) to
\(\langle q_j,w,0,0\rangle\).  The desired change to the proof
compiler is therefore exactly
\[
 \boxed{\quad
 \Gamma(B):\qquad\mathsf{select}_{1}(B,0,0)
       =\mathsf{Discard};\mathsf{AX}^{D}_{0,B,0};S^2;\mathsf{Pop}^{8}.
 \quad}
 \tag{D-Sel}\label{eq:gammad-selector}
\]
After Discard, \(d=h=0\), so this AX has \(D=B\).  The former
selector used \(\mathsf{select}_1(0,0,B)\).  Selection is safe even
when \(B=0\) and even when the unrelated gate guard fails.
The six retained heads use \(j=2,\ldots,7\), with their original
parameters and words.

For completeness, let \(\mathsf{Offset}_0\) be empty and let
\(\mathsf{Offset}_u=\mathsf{MP}_u\) for \(u>0\); from
\(\langle0,w,0,0\rangle\) it produces
\(\langle0,w,u,0\rangle\).  For \(u\le a\), the word
\[
 \mathsf{Discard};\mathsf{Offset}_u;
 \mathsf{AX}^{D}_{b,a-u,c};S^9;\mathsf{Pop}
\]
selects \(\rho_8(a,b,c,u)\).  The word
\(\mathsf{Discard};\mathsf{AX}^{D}_{c,b,a};\mathsf{Pop}^{10}\)
selects \(K(a,b,c)\).  Finally
\[
 \mathsf{Gen}_i=\mathsf{Push};\mathsf{Offset}_i;S;
                                      \mathsf{AX}^{D}_{i,0,0};S^{10}
\]
sends \(\langle a,w,0,0\rangle\) to
\(\langle U(i,a),w,0,0\rangle\), also when \(a=1\).
The Shift before AX tests zero and puts \(a\) in the cache; AX
clears that cache before a later Shift can test it.

\subsection{A strong proof compiler and both logical directions}

The cached-MP suffix is
\[
 \mathsf{MPSuffix}_b=S;S;\mathsf{MP}_b;S.
\]
From \(\langle0,[a,I(a,b),0;w],0,0\rangle\), it halts if
\(a=1\) or \(b=1\), and otherwise returns
\(\langle b,w,0,0\rangle\).  The second Shift tests \(a\);
successful MP places \(b\) in the cache; the final Shift tests \(b\)
and retrieves the dummy zero.  The dummy is necessary to preserve
the underlying \(w\).  This calculation also holds at the numerical
zero-stack alias.

Compile axiom leaves with the selectors just given, and compile a
generalization node by appending \(\mathsf{Gen}_i\) to its premise
word.  If an MP node has minor proof \(p_0\) of \(a\) and major
proof \(p_1\) of \(I(a,b)\), use
\[
 W_p=\mathsf{Discard};\mathsf{Push};
     W_{p_1};\mathsf{Push};W_{p_0};\mathsf{Push};\mathsf{MPSuffix}_b.
\]

\begin{lemma}[Uniform Gamma-D proof compiler]
For every finite \(\mathcal C_s\)-proof \(p\) and every \(t,w\),
the word \(W_p\) from \(\langle t,w,0,0\rangle\) either halts
or returns \(\langle\operatorname{concl}(p),w,0,0\rangle\).
If \(p\) proves code one, then
\[
  \mathsf{AX}^{D}_{0,0,0};W_p;\mathsf{Pop}
\]
halts from every incoming four-component store.
\end{lemma}
\begin{proof}
Induct on the proof tree.  The selectors prove the axiom case, including
the changed selector (\ref{eq:gammad-selector}).  The generalization
contract proves that case.  In the MP case, Discard and Push introduce
a dummy zero; if neither recursive premise word halts, their contracts
and Push produce precisely the suffix's three-entry precondition.
The suffix gives either an early halt or the required conclusion with
the original tail.  An earlier halt remains a halt when a word is
extended.  For the second claim, the initial AX clears both caches,
so the first claim applies to its arbitrary resulting top and stack.
If no earlier halt occurs, the last Pop tests the derived top one
before inspecting the arbitrary underlying stack.
\end{proof}

This all-store clause matters because the inherited exhaustive
command-word scheduler retains its stores between trials.  A refuting
proof yields a finite universally halting word, so the scheduler either
halts earlier or halts when it reaches this word.  The inherited
word encoding and scheduler proof are unchanged.

Conversely, annotate \(t\), \(h\), and each decoded stack entry
with a \(\mathcal C_s\)-derivation.  Code zero has the fixed proof
\(\rho_8(0,0,0,0)=0\), which supplies all zero padding.
AX preserves these annotations: the six heads and a successful gate
are axioms; a failed gate is zero; \(\Gamma(D)\) is an unconditional
instance of the existing diagonal family; \(d\le d+x\) permits
the propositional head; and \(U(d,h)\) follows by generalization
from the old cache.  A successful MP derives \(d+x\) from the
annotations for \(t\) and \(h=I(t,d+x)\); a failed MP returns
zero.  Shift moves derived entries, and a halting Shift exposes a
proof of one.  No derivability condition is required of the offset
\(d\).  The actual seeded source begins with the required zero
annotations.  This proves that its halting yields a refutation.

All selectors and proof-word constructors are primitive recursive.
The arguments are inductions on proof trees, finite lists, source loop
counters, and traces.  Together with the unchanged scheduler they prove
in PA that the ordinary Gamma-D cached-C source halts exactly when
\(\mathcal C_s\) has a refutation.  The proof translations in
(\ref{eq:gammad-cal}) then give its equivalence to
\(\neg\operatorname{Con}(\mathrm{PA})\).  In particular, the
soundness direction does not appeal to the standard truth of arithmetic
or to reflection for PA.

\subsection{Exact source binding and the two zero tests}

Let \(P_D\) be that ordinary cached-C Gamma-D source and let
\(P_D^Z\) be the source implemented by PA194D.  Both have 3,649
primitive sites, numbered from zero.  Relative to the old 190-token
Gamma-A emitter, the copy macros at token indices 121, 122, and 123
change from \texttt{copyA} to \texttt{copyD}.  Their source spans
are respectively 2258--2283, 2284--2309, and 2310--2335.
Each span has 26 primitive instructions.  The six altered register
arguments, and all their exact successors, are
\begin{center}\small
\begin{tabular}{@{}rll@{}}
\toprule
Site & Old instruction & Gamma-D instruction\\
\midrule
2277 & \(\operatorname{dec}(3,2281,2278)\) & \(\operatorname{dec}(4,2281,2278)\)\\
2282 & \(\operatorname{inc}(3,2283)\) & \(\operatorname{inc}(4,2283)\)\\
2303 & \(\operatorname{dec}(3,2307,2304)\) & \(\operatorname{dec}(4,2307,2304)\)\\
2308 & \(\operatorname{inc}(3,2309)\) & \(\operatorname{inc}(4,2309)\)\\
2329 & \(\operatorname{dec}(3,2333,2330)\) & \(\operatorname{dec}(4,2333,2330)\)\\
2334 & \(\operatorname{inc}(3,2335)\) & \(\operatorname{inc}(4,2335)\)\\
\bottomrule
\end{tabular}
\end{center}
Here registers 3 and 4 are logical \(A\) and \(D\), respectively.
Every other instruction of \(P_D\), including every successor,
agrees with the old cached-C source.  A separate free-expression
calculation checks the entire stack and gate of
(\ref{eq:gammad-ax}) for symbolic incoming parameters.

The physical search reassociates the emitter tree.  The winning
grammar's complete resolved instruction graph equals that of an
independently balanced tree over the same macro word.  It also equals
the graph obtained by changing only these macros in the supplied
accepted association.  Thus the binding checks structured jump
destinations, not only equality of the flattened token sequence.

To pass from \(P_D\) to \(P_D^Z\), change precisely
\[
\begin{array}{c|cc}
\text{site}&P_D&P_D^Z\\ \hline
26&\operatorname{jump}(34)&\operatorname{dec}(6,34,27)\\
3553&\operatorname{jump}(3637)&\operatorname{dec}(4,3637,3554).
\end{array}
\]
The first tests scratch \(X\); the second is the final AX return
and tests \(D\).  This adds two branch outcomes: \(P_D\) has
4,749, and \(P_D^Z\) has 4,751.

Use the positivity abstraction of Section~\ref{sec:invariant}, now
computed for \(P_D\) from the same clean seed
\[
 c_*=(25,(2,0,0,0,0,0,0,0,0)).
\]
Its least abstract closure \(R_D\) contains 263,244 location/mask
pairs.  A separate replay verifies that the seed belongs to it,
that it is closed, that the two source programs have exactly the same
abstract successors on it, and that the tested bits are zero at the
two changed sites.  There are 32 masks at site 26, all with bit 6
clear, and eight masks at site 3553, all with bit 4 clear.
The sorted-pair certificate has SHA-256
\begin{quote}\small
\nolinkurl{57c58dcae511fe13a6c200e7cfe4f4322370aec72b4a2fe5ba39b906fe62dbf8}.
\end{quote}
This is a fresh invariant for the changed formula emitter; it is not
an assumption that the Gamma-A invariant still applies.

The abstraction lemma and induction on execution length therefore
prove exact equality of the \(P_D\) and \(P_D^Z\) source runs
from \(c_*\), until and including any halt.  Each inserted decrement
sees zero, preserves every register, and takes the former jump target.
PA proves the finite closure facts and the two inductions, hence
\[
 \mathrm{PA}\vdash
 \bigl(\operatorname{Halt}(P_D,c_*)
          \leftrightarrow\operatorname{Halt}(P_D^Z,c_*)\bigr).
\]
This source equivalence is claimed from the specified seed.  The
arbitrary-store command calculations above concern \(P_D\);
they are not being inferred for arbitrary inputs to the specialized
source \(P_D^Z\).

\subsection{Literal realization and the PA theorem}

The nine folds are performed simultaneously on the 203-control raw
allocation. In native labels their replacement rows are:
\begin{center}\small
\begin{tabular}{@{}rrll@{}}
\toprule
Removed & Retained & New read-0 action & New read-1 action\\
\midrule
2004 & 2063 & $0\,L\,2049$ & $1\,R\,90$\\
2015 & 2023 & $0\,L\,2061$ & $1\,R\,2012$\\
2018 & 2012 & $0\,L\,2055$ & $1\,R\,90$\\
2022 & 2033 & $0\,L\,2002$ & $1\,R\,61$\\
2025 & 2046 & $0\,L\,2034$ & $1\,R\,2029$\\
2028 & 2039 & $0\,L\,2014$ & $1\,L\,2003$\\
2037 & 2059 & $0\,L\,2044$ & $1\,R\,2038$\\
2053 & 2011 & $0\,R\,98$ & $1\,L\,2019$\\
2064 & 2032 & $0\,L\,2017$ & $1\,L\,2041$\\
\bottomrule
\end{tabular}
\end{center}
Every incoming row establishes the required predecessor bit. None of the
folded controls is the initial control. The 36 local cases are checked
on the final simultaneous quotient: 14 take one step, 19 take three,
and three take five. The traces preserve the exterior tape and the full
required local effect; the supplied universal source contracts also
verify the resulting whole interpreter independently.

The literal table has an independently checked physical simulation
contract for each of the 4,751 primitive outcomes of \(P_D^Z\), including the positive
branches that the seeded source invariant excludes.  These contracts
range over arbitrary natural primary values and header gap and every
finite nonempty spare list.  The independent checker uses 1,091,154
symbolic steps, with a maximum of 385 for one case; every duration has
nonnegative coefficients and constant at least three.  It derives four
list lemmas from the literal rows and invokes no arithmetic-box oracle.
The nine compatible reader folds have 36 local all-tape cases.
The exact blank initialization takes 58,168
literal transitions to the complete frame for \(c_*\), in native
control 84 with head at \(-2\).  The frame has source path 1000000,
header distance 48, primary vector \((2,0,0,0,0,0,0,0,0)\), and
26 spare blocks of value one.

The source simulations terminate in positive time.  Consequently a
source halt produces a physical halt, and a physical halt cannot
precede a nonhalting concatenation of more positive source simulations
than the alleged halting time.  As in Section~\ref{sec:physical}, this
argument requires no inverse decoding of arbitrary intermediate tapes.
The finite startup and local certificates, and the unary, list and
trace inductions, are formalizable in PA.

\begin{theorem}[Gamma-D consistency machine]
Let \(H_D\) be the arithmetic sentence asserting a finite halting
trace of the specified PA194D table from its blank start.  Then
\[
 \mathrm{PA}\vdash
    \bigl(H_D\leftrightarrow\neg\operatorname{Con}(\mathrm{PA})\bigr),
 \qquad
 \mathrm{PA}\vdash
    \bigl(\neg H_D\leftrightarrow\operatorname{Con}(\mathrm{PA})\bigr).
\]
\end{theorem}
\begin{proof}
Compose, inside PA, the physical simulation, the seeded zero-test
equivalence, the soundness and strong proof compiler, and the supplied
proof calibration:
\[
\begin{split}
 H_D
 &\longleftrightarrow\operatorname{Halt}(P_D^Z,c_*)
 \longleftrightarrow\operatorname{Halt}(P_D,c_*)\\
 &\longleftrightarrow
       \exists q\,\operatorname{Prf}_{\mathcal C_s}(q,1)
 \longleftrightarrow\neg\operatorname{Con}(\mathrm{PA}).
\end{split}
\]
Classical logic in PA gives the nonhalting formulation.
\end{proof}

The theorem is an ordinary proof with explicit finite and symbolic
certificates, using the retained ordinary PA calibration.  It is not
a new proof-assistant kernel export or an iPRA equivalence theorem.
It makes no consistency assumption in establishing the PA-provable
biconditional.  PA194D supplies a second construction at the requested
194-state fallback threshold; it does not lower that threshold.
